% !TEX program = xelatex
\documentclass[9pt,aspectratio=169]{beamer}
\usetheme{metropolis}

% ── Packages ──────────────────────────────────────────────
\usepackage{amsmath,amssymb}
\usepackage{booktabs}
\usepackage{graphicx}
\usepackage{tikz}
\usepackage{appendixnumberbeamer}

% ── Meta ──────────────────────────────────────────────────
\title{Bundling, Product Mix, and SME Outcomes\\in Public Procurement}
\subtitle{Cross-Sectional Evidence from Chile's Compra Ágil}
\author{Nicolás Jiménez}
\date{April 2026}
\institute{Yale University}

\begin{document}

% ══════════════════════════════════════════════════════════
% TITLE
% ══════════════════════════════════════════════════════════
\maketitle

% ══════════════════════════════════════════════════════════
\section{Motivation \& Setup}
% ══════════════════════════════════════════════════════════

\begin{frame}{Motivation}
After the Compra Ágil reform, municipal buyers \textbf{bundle more products} into each tender:
\begin{itemize}
  \item More line items per tender (+0.64)
  \item More distinct 8-digit product codes (+0.28)
  \item More 2-digit UNSPSC segments (+0.06)
\end{itemize}

\bigskip
\textbf{Question:} Does this bundling help or hurt SMEs and local firms?

\bigskip
Two dimensions of bundle breadth may have opposite effects:
\begin{itemize}
  \item \textbf{Product variety} (N distinct product codes overall) $\to$ favors large firms with broad catalogs?
  \item \textbf{Category breadth} (N segments across categories) $\to$ favors local generalist suppliers?
\end{itemize}
\end{frame}

\begin{frame}{Related Literature}
\textbf{Bundling in procurement:} fundamental tradeoff between production efficiencies and competition effects.
\begin{itemize}
  \item \textbf{Scale economies vs.\ entry deterrence:}
        Bundling can reduce unit costs but weaken competition through entry barriers
        (Estache \& Iimi 2011; Qiao et al.\ 2019; Ridderstedt \& Nilsson 2023)
  \item \textbf{Bundling as multi-dimensional treatment:}
        Swedish pavement studies decompose into scale, spatial spread, task heterogeneity ---
        net effect depends on \textit{what} is bundled, not just \textit{how much}
        (Ridderstedt \& Nilsson 2023)
  \item \textbf{Bidder asymmetry:}
        Bundling raises incumbent advantage and excludes entrants
        (Suzuki 2021, electricity; Lunander \& Lundberg 2012, roads)
  \item \textbf{Monitoring \& quality:}
        Unbundling + inspections improve quality but delay delivery;
        monitoring can substitute for unbundling
        (Wolfram et al.\ 2023, Kenya electrification)
\end{itemize}

\medskip
\textbf{Gap we fill:} Most studies treat bundling as a scalar. We decompose it into \textit{product variety} vs.\ \textit{category breadth} and separate \textit{selection} (who bids) from \textit{competitive advantage} (who wins conditional on bidding).
\end{frame}

\begin{frame}{Empirical Approach}
\textbf{Sample:} Municipal Compra Ágil tenders, Jan 2023--Mar 2025.

\bigskip
\textbf{Specification:}
\[
  y_{it} = \alpha_i + \gamma_t + \beta_1 \, \text{N\_product8}_{it} + \beta_2 \, \text{N\_segment2}_{it} + \mathbf{X}_{it}'\delta + \varepsilon_{it}
\]

\begin{itemize}
  \item $\alpha_i$: buyer fixed effects (or region FE); $\gamma_t$: year-month FE
  \item \textbf{N\_product8}: distinct 8-digit UNSPSC commodity codes in the tender (product variety)
  \item \textbf{N\_segment2}: distinct 2-digit UNSPSC segments in the tender (category breadth)
  \item Controls: log(estimated cost), N bidders, top-20 segment dummies
  \item SEs clustered by buyer
  \item Separate estimates for pre-reform and post-reform periods
\end{itemize}
\end{frame}

% ══════════════════════════════════════════════════════════
\section{Who Wins? Baseline Results}
% ══════════════════════════════════════════════════════════


\begin{frame}{Pr(SME Wins) on Bundle Composition}
\centering\small
\begin{tabular}{lcccc}
\toprule
 & Pre & Pre + Seg & Post & Post + Seg \\
 & (1) & (2) & (3) & (4) \\
\midrule
N product codes (8-digit) & -0.0091$^{***}$ & -0.0015$^{**}$ & -0.0045$^{***}$ & -0.0013$^{**}$ \\
 & (0.0010) & (0.0006) & (0.0009) & (0.0005) \\[2pt]
N segments (2-digit) & 0.0263$^{***}$ & 0.0199$^{***}$ & 0.0178$^{***}$ & 0.0072$^{***}$ \\
 & (0.0025) & (0.0022) & (0.0024) & (0.0017) \\[2pt]
log(Estimated cost) & -0.0013 & 0.0032$^{**}$ & -0.0104$^{***}$ & -0.0065$^{***}$ \\
 & (0.0022) & (0.0016) & (0.0010) & (0.0008) \\[2pt]
N bidders & -0.0051$^{***}$ & -0.0036$^{***}$ & 0.0046$^{***}$ & 0.0043$^{***}$ \\
 & (0.0003) & (0.0003) & (0.0003) & (0.0002) \\[2pt]
\midrule
Buyer + YM FE & Yes & Yes & Yes & Yes \\
Segment dummies & No & Yes & No & Yes \\
N & 545,518 & 545,518 & 155,434 & 155,434 \\
Within $R^2$ & 0.005 & 0.090 & 0.014 & 0.102 \\
\bottomrule
\multicolumn{5}{l}{\footnotesize SEs clustered by buyer. $^{***}p<0.01$, $^{**}p<0.05$, $^{*}p<0.10$.} \\
\end{tabular}

\medskip
\small\textbf{Takeaway:} Product variety hurts SMEs; category breadth helps them. Segment dummies absorb much of the raw effect.
\end{frame}

\begin{frame}{Pr(Local SME Wins) on Bundle Composition}
\centering\small
\begin{tabular}{lcccc}
\toprule
 & Pre & Pre + Seg & Post & Post + Seg \\
 & (1) & (2) & (3) & (4) \\
\midrule
N product codes (8-digit) & -0.0011 & 0.0001 & 0.0004 & -0.0007 \\
 & (0.0010) & (0.0007) & (0.0013) & (0.0008) \\[2pt]
N segments (2-digit) & 0.0101$^{***}$ & 0.0605$^{***}$ & 0.0003 & 0.0594$^{***}$ \\
 & (0.0023) & (0.0030) & (0.0032) & (0.0037) \\[2pt]
log(Estimated cost) & 0.0034 & 0.0050$^{**}$ & -0.0049$^{**}$ & -0.0042$^{**}$ \\
 & (0.0025) & (0.0021) & (0.0022) & (0.0019) \\[2pt]
N bidders & -0.0165$^{***}$ & -0.0108$^{***}$ & -0.0128$^{***}$ & -0.0078$^{***}$ \\
 & (0.0007) & (0.0005) & (0.0009) & (0.0007) \\[2pt]
\midrule
Buyer + YM FE & Yes & Yes & Yes & Yes \\
Segment dummies & No & Yes & No & Yes \\
N & 545,518 & 545,518 & 155,433 & 155,433 \\
Within $R^2$ & 0.028 & 0.094 & 0.020 & 0.105 \\
\bottomrule
\multicolumn{5}{l}{\footnotesize SEs clustered by buyer. $^{***}p<0.01$, $^{**}p<0.05$, $^{*}p<0.10$.} \\
\end{tabular}

\medskip
\small\textbf{Takeaway:} N\_segment2 has a strong positive effect on local SME wins (+0.06$^{***}$) once segment composition is controlled.
\end{frame}

% ══════════════════════════════════════════════════════════
\section{Selection vs.\ Competitive Advantage}
% ══════════════════════════════════════════════════════════

\begin{frame}{Is It Selection or Competitive Advantage?}
The baseline results could reflect two channels:
\begin{enumerate}
  \item \textbf{Selection:} Bundle composition changes \textit{who bids} \\
        $\to$ More product variety attracts large firms into the bidder pool
  \item \textbf{Competitive advantage:} Bundle composition changes \textit{who wins conditional on bidding} \\
        $\to$ Cross-category breadth gives local SMEs an edge
\end{enumerate}

\bigskip
\textbf{Test:} Add \textbf{share of SME bidders} as a control.
\begin{itemize}
  \item If results survive: competitive advantage channel
  \item If results disappear: selection channel
\end{itemize}
\end{frame}

\begin{frame}{Pr(SME Wins): Adding SME Share of Bidders}
\centering\small
\begin{tabular}{lcccc}
\toprule
 & Pre & Pre + SME\% & Post & Post + SME\% \\
 & (1) & (2) & (3) & (4) \\
\midrule
N product codes (8-digit) & -0.0015$^{**}$ & 0.0052$^{***}$ & -0.0013$^{**}$ & -0.0010$^{***}$ \\
 & (0.0006) & (0.0005) & (0.0005) & (0.0002) \\[2pt]
N segments (2-digit) & 0.0199$^{***}$ & 0.0008 & 0.0072$^{***}$ & 0.0043$^{***}$ \\
 & (0.0022) & (0.0011) & (0.0017) & (0.0008) \\[2pt]
Share SME bidders &  & 1.1571$^{***}$ &  & 1.1107$^{***}$ \\
 &  & (0.0035) &  & (0.0052) \\[2pt]
log(Estimated cost) & 0.0032$^{**}$ & 0.0057$^{***}$ & -0.0065$^{***}$ & 0.0003 \\
 & (0.0016) & (0.0007) & (0.0008) & (0.0004) \\[2pt]
N bidders & -0.0036$^{***}$ & -0.0046$^{***}$ & 0.0043$^{***}$ & -0.0004$^{***}$ \\
 & (0.0003) & (0.0002) & (0.0002) & (0.0001) \\[2pt]
\midrule
Buyer + YM FE & Yes & Yes & Yes & Yes \\
Segment dummies & Yes & Yes & Yes & Yes \\
SME share control & No & Yes & No & Yes \\
N & 545,518 & 545,518 & 155,434 & 155,434 \\
Within $R^2$ & 0.090 & 0.495 & 0.102 & 0.680 \\
\bottomrule
\multicolumn{5}{l}{\footnotesize SEs clustered by buyer. $^{***}p<0.01$, $^{**}p<0.05$, $^{*}p<0.10$.} \\
\end{tabular}

\medskip
\small\textbf{Takeaway:} SME share absorbs almost all variation ($R^2$ jumps to $\sim$0.49). Product variety effect flips sign --- selection, not competitive disadvantage.
\end{frame}

\begin{frame}{Pr(Local SME Wins): Adding SME Share of Bidders}
\centering\small
\begin{tabular}{lcccc}
\toprule
 & Pre & Pre + SME\% & Post & Post + SME\% \\
 & (1) & (2) & (3) & (4) \\
\midrule
N product codes (8-digit) & 0.0001 & 0.0039$^{***}$ & -0.0007 & -0.0006 \\
 & (0.0007) & (0.0007) & (0.0008) & (0.0008) \\[2pt]
N segments (2-digit) & 0.0605$^{***}$ & 0.0496$^{***}$ & 0.0594$^{***}$ & 0.0577$^{***}$ \\
 & (0.0030) & (0.0029) & (0.0037) & (0.0035) \\[2pt]
Share SME bidders &  & 0.6594$^{***}$ &  & 0.6317$^{***}$ \\
 &  & (0.0148) &  & (0.0173) \\[2pt]
log(Estimated cost) & 0.0050$^{**}$ & 0.0065$^{***}$ & -0.0042$^{**}$ & -0.0003 \\
 & (0.0021) & (0.0018) & (0.0019) & (0.0018) \\[2pt]
N bidders & -0.0108$^{***}$ & -0.0113$^{***}$ & -0.0078$^{***}$ & -0.0105$^{***}$ \\
 & (0.0005) & (0.0005) & (0.0007) & (0.0006) \\[2pt]
\midrule
Buyer + YM FE & Yes & Yes & Yes & Yes \\
Segment dummies & Yes & Yes & Yes & Yes \\
SME share control & No & Yes & No & Yes \\
N & 545,518 & 545,518 & 155,433 & 155,433 \\
Within $R^2$ & 0.094 & 0.187 & 0.105 & 0.151 \\
\bottomrule
\multicolumn{5}{l}{\footnotesize SEs clustered by buyer. $^{***}p<0.01$, $^{**}p<0.05$, $^{*}p<0.10$.} \\
\end{tabular}

\medskip
\small\textbf{Key result:} N\_segment2 effect \textbf{survives} at $+0.050^{***}$ (pre) and $+0.058^{***}$ (post). \\
Cross-category breadth helps local SMEs win \textit{conditional on who bids} --- genuine competitive advantage.
\end{frame}

\begin{frame}{Pr(Large Firm Wins): Adding SME Share of Bidders}
\centering\small
\begin{tabular}{lcccc}
\toprule
 & Pre & Pre + SME\% & Post & Post + SME\% \\
 & (1) & (2) & (3) & (4) \\
\midrule
N product codes (8-digit) & 0.0015$^{**}$ & -0.0052$^{***}$ & 0.0013$^{**}$ & 0.0010$^{***}$ \\
 & (0.0006) & (0.0005) & (0.0005) & (0.0002) \\[2pt]
N segments (2-digit) & -0.0199$^{***}$ & -0.0008 & -0.0072$^{***}$ & -0.0043$^{***}$ \\
 & (0.0022) & (0.0011) & (0.0017) & (0.0008) \\[2pt]
Share SME bidders &  & -1.1571$^{***}$ &  & -1.1107$^{***}$ \\
 &  & (0.0035) &  & (0.0052) \\[2pt]
log(Estimated cost) & -0.0032$^{**}$ & -0.0057$^{***}$ & 0.0065$^{***}$ & -0.0003 \\
 & (0.0016) & (0.0007) & (0.0008) & (0.0004) \\[2pt]
N bidders & 0.0036$^{***}$ & 0.0046$^{***}$ & -0.0043$^{***}$ & 0.0004$^{***}$ \\
 & (0.0003) & (0.0002) & (0.0002) & (0.0001) \\[2pt]
\midrule
Buyer + YM FE & Yes & Yes & Yes & Yes \\
Segment dummies & Yes & Yes & Yes & Yes \\
SME share control & No & Yes & No & Yes \\
N & 545,518 & 545,518 & 155,434 & 155,434 \\
Within $R^2$ & 0.090 & 0.495 & 0.102 & 0.680 \\
\bottomrule
\multicolumn{5}{l}{\footnotesize SEs clustered by buyer. $^{***}p<0.01$, $^{**}p<0.05$, $^{*}p<0.10$.} \\
\end{tabular}

\medskip
\small Mirror image of SME results. SME share of bidders is the dominant predictor.
\end{frame}

% ══════════════════════════════════════════════════════════
\section{The Selection Channel: Who Bids?}
% ══════════════════════════════════════════════════════════

\begin{frame}{Bidder Composition as Outcome}
If bundling changes \textit{who bids}, we should see it directly in the share of SME and local bidders.

\bigskip
\textbf{Outcome variables:}
\begin{itemize}
  \item Share of bidders that are SMEs
  \item Share of bidders that are local SMEs
  \item Share of bidders that are large firms
\end{itemize}

\medskip
Same controls as before (minus N bidders, which is mechanically related), plus segment dummies.
\end{frame}

\begin{frame}{Share of SME Bidders on Bundle Composition}
\centering\small
\begin{tabular}{lcccc}
\toprule
 & Pre & Pre + Seg & Post & Post + Seg \\
 & (1) & (2) & (3) & (4) \\
\midrule
N product codes (8-digit) & -0.0120$^{***}$ & -0.0061$^{***}$ & -0.0026$^{***}$ & -0.0003 \\
 & (0.0008) & (0.0005) & (0.0005) & (0.0003) \\[2pt]
N segments (2-digit) & 0.0157$^{***}$ & 0.0195$^{***}$ & 0.0103$^{***}$ & 0.0045$^{***}$ \\
 & (0.0017) & (0.0014) & (0.0013) & (0.0009) \\[2pt]
log(Estimated cost) & -0.0036$^{**}$ & -0.0019$^{**}$ & -0.0059$^{***}$ & -0.0038$^{***}$ \\
 & (0.0014) & (0.0009) & (0.0006) & (0.0005) \\[2pt]
N bidders & -0.0024$^{***}$ & -0.0007$^{***}$ & 0.0031$^{***}$ & 0.0030$^{***}$ \\
 & (0.0002) & (0.0001) & (0.0002) & (0.0001) \\[2pt]
\midrule
Buyer + YM FE & Yes & Yes & Yes & Yes \\
Segment dummies & No & Yes & No & Yes \\
N & 880,029 & 880,029 & 288,525 & 288,525 \\
Within $R^2$ & 0.012 & 0.158 & 0.014 & 0.104 \\
\bottomrule
\multicolumn{5}{l}{\footnotesize SEs clustered by buyer. $^{***}p<0.01$, $^{**}p<0.05$, $^{*}p<0.10$.} \\
\end{tabular}

\medskip
\small Product variety reduces SME share of bidders ($-0.006^{***}$). Category breadth increases it ($+0.020^{***}$).
\end{frame}

\begin{frame}{Share of Local SME Bidders on Bundle Composition}
\centering\small
\begin{tabular}{lcccc}
\toprule
 & Pre & Pre + Seg & Post & Post + Seg \\
 & (1) & (2) & (3) & (4) \\
\midrule
N product codes (8-digit) & -0.0013$^{*}$ & -0.0019$^{***}$ & 0.0026$^{***}$ & 0.0007 \\
 & (0.0008) & (0.0005) & (0.0009) & (0.0005) \\[2pt]
N segments (2-digit) & 0.0046$^{***}$ & 0.0398$^{***}$ & -0.0029 & 0.0343$^{***}$ \\
 & (0.0016) & (0.0018) & (0.0021) & (0.0022) \\[2pt]
log(Estimated cost) & -0.0072$^{***}$ & -0.0050$^{***}$ & -0.0055$^{***}$ & -0.0038$^{***}$ \\
 & (0.0016) & (0.0013) & (0.0014) & (0.0012) \\[2pt]
N bidders & -0.0096$^{***}$ & -0.0051$^{***}$ & -0.0084$^{***}$ & -0.0052$^{***}$ \\
 & (0.0004) & (0.0003) & (0.0005) & (0.0004) \\[2pt]
\midrule
Buyer + YM FE & Yes & Yes & Yes & Yes \\
Segment dummies & No & Yes & No & Yes \\
N & 880,029 & 880,029 & 288,525 & 288,525 \\
Within $R^2$ & 0.019 & 0.110 & 0.018 & 0.109 \\
\bottomrule
\multicolumn{5}{l}{\footnotesize SEs clustered by buyer. $^{***}p<0.01$, $^{**}p<0.05$, $^{*}p<0.10$.} \\
\end{tabular}

\medskip
\small Strong positive effect of N\_segment2 on local SME share ($+0.040^{***}$ pre, $+0.034^{***}$ post). \\
Cross-category bundling shifts the bidder pool toward local SMEs.
\end{frame}

\begin{frame}{Share of Large-Firm Bidders on Bundle Composition}
\centering\small
\begin{tabular}{lcccc}
\toprule
 & Pre & Pre + Seg & Post & Post + Seg \\
 & (1) & (2) & (3) & (4) \\
\midrule
N product codes (8-digit) & 0.0112$^{***}$ & 0.0057$^{***}$ & 0.0025$^{***}$ & 0.0003 \\
 & (0.0008) & (0.0005) & (0.0005) & (0.0003) \\[2pt]
N segments (2-digit) & -0.0157$^{***}$ & -0.0191$^{***}$ & -0.0102$^{***}$ & -0.0044$^{***}$ \\
 & (0.0017) & (0.0014) & (0.0013) & (0.0009) \\[2pt]
log(Estimated cost) & 0.0040$^{***}$ & 0.0019$^{**}$ & 0.0059$^{***}$ & 0.0038$^{***}$ \\
 & (0.0014) & (0.0009) & (0.0006) & (0.0005) \\[2pt]
N bidders & 0.0021$^{***}$ & 0.0004$^{***}$ & -0.0031$^{***}$ & -0.0030$^{***}$ \\
 & (0.0002) & (0.0001) & (0.0002) & (0.0001) \\[2pt]
\midrule
Buyer + YM FE & Yes & Yes & Yes & Yes \\
Segment dummies & No & Yes & No & Yes \\
N & 880,029 & 880,029 & 288,525 & 288,525 \\
Within $R^2$ & 0.011 & 0.168 & 0.015 & 0.106 \\
\bottomrule
\multicolumn{5}{l}{\footnotesize SEs clustered by buyer. $^{***}p<0.01$, $^{**}p<0.05$, $^{*}p<0.10$.} \\
\end{tabular}

\medskip
\small Mirror image: product variety attracts large firms ($+0.006^{***}$); category breadth repels them ($-0.019^{***}$).
\end{frame}

% ══════════════════════════════════════════════════════════
\section{Extensive Margin: Number of Bidders}
% ══════════════════════════════════════════════════════════

\begin{frame}{How Does Bundling Affect Participation?}
The share results show bundling tilts the composition of bidders. But does it also change the \textbf{total number} of bidders?

\bigskip
\textbf{Outcome variables:}
\begin{itemize}
  \item N total bidders
  \item N SME bidders
  \item N large-firm bidders
  \item N local bidders
\end{itemize}

\medskip
Controls: log(estimated cost), segment dummies, buyer + year-month FE. \\
\textit{N bidders is now an outcome, not a control.}
\end{frame}

\begin{frame}{N Total Bidders on Bundle Composition}
\centering\small
\begin{tabular}{lcccc}
\toprule
 & Pre & Pre + Seg & Post & Post + Seg \\
 & (1) & (2) & (3) & (4) \\
\midrule
N product codes (8-digit) & 0.0014 & 0.0069 & -0.0299$^{***}$ & -0.0061 \\
 & (0.0076) & (0.0065) & (0.0063) & (0.0060) \\[2pt]
N segments (2-digit) & -0.0721$^{***}$ & -1.0823$^{***}$ & -0.0019 & -1.2045$^{***}$ \\
 & (0.0168) & (0.0236) & (0.0161) & (0.0300) \\[2pt]
log(Estimated cost) & 0.2911$^{***}$ & 0.3180$^{***}$ & 0.2251$^{***}$ & 0.2992$^{***}$ \\
 & (0.0169) & (0.0134) & (0.0173) & (0.0151) \\[2pt]
\midrule
Buyer + YM FE & Yes & Yes & Yes & Yes \\
Segment dummies & No & Yes & No & Yes \\
N & 930,247 & 930,247 & 307,298 & 307,298 \\
Within $R^2$ & 0.003 & 0.120 & 0.003 & 0.111 \\
\bottomrule
\multicolumn{5}{l}{\footnotesize SEs clustered by buyer. $^{***}p<0.01$, $^{**}p<0.05$, $^{*}p<0.10$.} \\
\end{tabular}

\medskip
\small With segment dummies, N\_segment2 has a \textbf{large negative} effect ($-1.08^{***}$): each additional segment \textit{reduces} total bidders by $\sim$1. Diverse tenders are harder to bid on.
\end{frame}

\begin{frame}{N SME Bidders on Bundle Composition}
\centering\small
\begin{tabular}{lcccc}
\toprule
 & Pre & Pre + Seg & Post & Post + Seg \\
 & (1) & (2) & (3) & (4) \\
\midrule
N product codes (8-digit) & -0.0614$^{***}$ & -0.0243$^{***}$ & -0.0405$^{***}$ & -0.0056 \\
 & (0.0050) & (0.0044) & (0.0062) & (0.0058) \\[2pt]
N segments (2-digit) & 0.0558$^{***}$ & -0.8513$^{***}$ & 0.0468$^{***}$ & -1.1287$^{***}$ \\
 & (0.0116) & (0.0177) & (0.0153) & (0.0274) \\[2pt]
log(Estimated cost) & 0.1987$^{***}$ & 0.2310$^{***}$ & 0.1896$^{***}$ & 0.2524$^{***}$ \\
 & (0.0149) & (0.0113) & (0.0152) & (0.0130) \\[2pt]
\midrule
Buyer + YM FE & Yes & Yes & Yes & Yes \\
Segment dummies & No & Yes & No & Yes \\
N & 930,247 & 930,247 & 307,298 & 307,298 \\
Within $R^2$ & 0.003 & 0.125 & 0.002 & 0.128 \\
\bottomrule
\multicolumn{5}{l}{\footnotesize SEs clustered by buyer. $^{***}p<0.01$, $^{**}p<0.05$, $^{*}p<0.10$.} \\
\end{tabular}

\medskip
\small Product variety reduces SME bidders ($-0.024^{***}$). Category breadth also reduces SME bidders with segment controls ($-0.85^{***}$) --- but recall it \textit{increases their share}.
\end{frame}

\begin{frame}{N Large-Firm Bidders on Bundle Composition}
\centering\small
\begin{tabular}{lcccc}
\toprule
 & Pre & Pre + Seg & Post & Post + Seg \\
 & (1) & (2) & (3) & (4) \\
\midrule
N product codes (8-digit) & 0.0656$^{***}$ & 0.0363$^{***}$ & 0.0116$^{***}$ & 0.0026$^{**}$ \\
 & (0.0059) & (0.0036) & (0.0019) & (0.0010) \\[2pt]
N segments (2-digit) & -0.1060$^{***}$ & -0.2622$^{***}$ & -0.0383$^{***}$ & -0.0416$^{***}$ \\
 & (0.0129) & (0.0093) & (0.0046) & (0.0029) \\[2pt]
log(Estimated cost) & 0.0840$^{***}$ & 0.0669$^{***}$ & 0.0188$^{***}$ & 0.0107$^{***}$ \\
 & (0.0079) & (0.0043) & (0.0019) & (0.0012) \\[2pt]
\midrule
Buyer + YM FE & Yes & Yes & Yes & Yes \\
Segment dummies & No & Yes & No & Yes \\
N & 930,247 & 930,247 & 307,298 & 307,298 \\
Within $R^2$ & 0.012 & 0.230 & 0.006 & 0.205 \\
\bottomrule
\multicolumn{5}{l}{\footnotesize SEs clustered by buyer. $^{***}p<0.01$, $^{**}p<0.05$, $^{*}p<0.10$.} \\
\end{tabular}

\medskip
\small Product variety attracts large firms ($+0.036^{***}$). Category breadth repels them ($-0.26^{***}$). \\
The reduction in large-firm participation from cross-category bundling is proportionally larger than for SMEs.
\end{frame}

\begin{frame}{N Local Bidders on Bundle Composition}
\centering\small
\begin{tabular}{lcccc}
\toprule
 & Pre & Pre + Seg & Post & Post + Seg \\
 & (1) & (2) & (3) & (4) \\
\midrule
N product codes (8-digit) & -0.0154$^{***}$ & -0.0182$^{***}$ & -0.0110$^{**}$ & -0.0108$^{***}$ \\
 & (0.0045) & (0.0038) & (0.0046) & (0.0039) \\[2pt]
N segments (2-digit) & 0.0600$^{***}$ & -0.2420$^{***}$ & 0.0376$^{***}$ & -0.2815$^{***}$ \\
 & (0.0085) & (0.0224) & (0.0084) & (0.0276) \\[2pt]
log(Estimated cost) & 0.0622$^{***}$ & 0.0921$^{***}$ & 0.0369$^{***}$ & 0.0725$^{***}$ \\
 & (0.0085) & (0.0082) & (0.0084) & (0.0085) \\[2pt]
\midrule
Buyer + YM FE & Yes & Yes & Yes & Yes \\
Segment dummies & No & Yes & No & Yes \\
N & 930,247 & 930,247 & 307,298 & 307,298 \\
Within $R^2$ & 0.001 & 0.036 & 0.000 & 0.034 \\
\bottomrule
\multicolumn{5}{l}{\footnotesize SEs clustered by buyer. $^{***}p<0.01$, $^{**}p<0.05$, $^{*}p<0.10$.} \\
\end{tabular}

\medskip
\small Category breadth reduces local bidders too ($-0.24^{***}$), but less than total ($-1.08^{***}$). \\
Diverse tenders shrink the pool but tilt it local --- consistent with the ``general store'' story.
\end{frame}

% ══════════════════════════════════════════════════════════
\section{Firm Heterogeneity: Generalists vs.\ Specialists}
% ══════════════════════════════════════════════════════════

\begin{frame}{Who Are the Generalists?}
\textbf{Firm scope} = N distinct 2-digit UNSPSC segments bid on in the pre-period (Jan 2022 -- Nov 2024). \\
\textbf{Generalist} = scope $\geq 10$ segments (top quartile of 67K firms).

\bigskip
\begin{columns}[T]
  \column{0.48\textwidth}
  \textbf{SII sector (rubro):}
  \begin{itemize}
    \item 63\% of generalists are in \textit{Comercio al por mayor y menor} (wholesale/retail trade)
    \item vs.\ 31\% of specialists
    \item Modal generalist is literally a \textbf{commercial trading firm} --- a distributor or general store
  \end{itemize}

  \bigskip
  \textbf{Firm size:}
  \begin{itemize}
    \item Generalists are modestly larger (mean tramo ventas 5.0 vs.\ 4.4; mean 26 vs.\ 18 workers)
    \item But correlation is low ($\rho = 0.12$) --- scope $\neq$ size
  \end{itemize}

  \column{0.48\textwidth}
  \textbf{Geography:}
  \begin{itemize}
    \item 42.6\% of generalists in Santiago vs.\ 37.1\% of specialists
    \item Modest difference --- generalists are spread nationwide
    \item Region HHI: 0.22 (gen) vs.\ 0.18 (spec)
  \end{itemize}

  \bigskip
  \textbf{Takeaway:}
  \begin{itemize}
    \item Generalists $\approx$ local/regional \textbf{trade \& distribution firms} with broad product catalogs
    \item Not simply ``big firms'' --- many are small
    \item Present across Chile, not just Santiago
  \end{itemize}
\end{columns}
\end{frame}

\begin{frame}{Who Are the Generalists? Size Distribution}
\centering\small
\begin{tabular}{llrr}
\toprule
Tramo & Sales bracket (UF) & Specialists & Generalists \\
\midrule
1 & $<$ 800 & 7.3\% & 4.4\% \\
2 & 800 -- 2,400 & 17.6\% & 11.2\% \\
3 & 2,400 -- 5,000 & 13.6\% & 11.5\% \\
4 & 5,000 -- 10,000 & 22.5\% & 23.3\% \\
5 & 10,000 -- 25,000 & 11.9\% & 13.6\% \\
6 & 25,000 -- 50,000 & 9.2\% & 10.7\% \\
7 & 50,000 -- 100,000 & 8.4\% & 10.9\% \\
8 & 100,000 -- 200,000 & 3.7\% & 5.5\% \\
9 & 200,000 -- 600,000 & 2.3\% & 3.1\% \\
10--13 & $>$ 600,000 & 3.5\% & 5.8\% \\
\midrule
\multicolumn{2}{l}{Mean tramo ventas} & 4.36 & 5.00 \\
\multicolumn{2}{l}{Mean workers} & 18.0 & 25.5 \\
\multicolumn{2}{l}{Median workers} & 0 & 1 \\
\multicolumn{2}{l}{N firms} & 32,982 & 6,732 \\
\bottomrule
\end{tabular}

\bigskip
\small Generalists shift rightward but overlap heavily with specialists. \\
Correlation of scope with tramo ventas = 0.12; with N workers = 0.01.
\end{frame}

\begin{frame}{Do Generalist Firms Benefit from Diverse Tenders?}
\textbf{Idea:} If cross-category bundling helps local generalists, we should see it
at the \textit{firm level} --- firms with broader pre-period scope should do better on diverse tenders.

\bigskip
\textbf{Setup:}
\begin{itemize}
  \item \textbf{Firm scope} = N distinct 2-digit UNSPSC segments a firm bid on in the pre-period (Jan 2022 -- Nov 2024)
  \item \textbf{Generalist} = firm scope $\geq 10$ segments (top quartile)
  \item Unit of observation: bid (6.6M bids across 1.2M tenders)
  \item Within-\textit{tender} FE $\to$ identifies from variation across bidders on the \textit{same tender}
\end{itemize}

\bigskip
\textbf{Key question:} Does the generalist penalty shrink after the reform induced more bundling?
\end{frame}

\begin{frame}{Generalist Penalty: Pre vs.\ Post Reform}
\centering\small
\begin{tabular}{lccc}
\toprule
 & Pre $\beta$ (SE) & Post $\beta$ (SE) & $\Delta$ \\
\midrule
\multicolumn{4}{l}{\textit{Panel A: Binary generalist ($\geq$10 segments)}} \\[3pt]
All tenders & $-$0.163 (0.001) & $-$0.065 (0.001) & +0.098$^{***}$ \\
1 segment tenders & $-$0.168 (0.001) & $-$0.073 (0.001) & +0.095$^{***}$ \\
2 segment tenders & $-$0.097 (0.004) & $-$0.027 (0.002) & +0.071$^{***}$ \\
3+ segment tenders & $-$0.090 (0.006) & $-$0.015 (0.003) & +0.075$^{***}$ \\
\midrule
\multicolumn{4}{l}{\textit{Panel B: Continuous log(1 + firm scope)}} \\[3pt]
All tenders & $-$0.109 (0.000) & $-$0.023 (0.000) & +0.086$^{***}$ \\
1 segment tenders & $-$0.109 (0.000) & $-$0.026 (0.000) & +0.083$^{***}$ \\
2 segment tenders & $-$0.092 (0.001) & $-$0.013 (0.001) & +0.079$^{***}$ \\
3+ segment tenders & $-$0.131 (0.002) & $-$0.009 (0.001) & +0.121$^{***}$ \\
\bottomrule
\multicolumn{4}{l}{\footnotesize Within-tender FE (bid-level). SEs clustered by tender.} \\
\multicolumn{4}{l}{\footnotesize Outcome: bid wins tender (0/1). Scope measured in pre-period only.} \\
\end{tabular}
\end{frame}

\begin{frame}{Interpretation: The General Store Advantage}
\begin{enumerate}
  \item \textbf{Pre-reform:} Generalist firms face a significant penalty ($-$16pp) when competing on the same tender as specialists
  \begin{itemize}
    \item Specialists have deeper expertise in narrow categories
    \item Broad scope is a liability in focused tenders
  \end{itemize}

  \bigskip
  \item \textbf{Post-reform:} The penalty shrinks by $\sim$60\% overall and nearly vanishes on diverse tenders
  \begin{itemize}
    \item On 3+ segment tenders: $-$9.0pp $\to$ $-$1.5pp (binary), $-$13.1pp $\to$ $-$0.9pp (continuous)
    \item Reform-induced bundling made cross-category capability valuable
  \end{itemize}

  \bigskip
  \item \textbf{Connects tender-level and firm-level stories:}
  \begin{itemize}
    \item Tender level: N\_segment2 $\uparrow$ $\to$ Pr(local SME wins) $\uparrow$
    \item Firm level: generalists gain ground on diverse tenders post-reform
    \item Mechanism: local generalist SMEs are the ``general stores'' that can fill cross-category orders
  \end{itemize}
\end{enumerate}
\end{frame}

\begin{frame}{SII Validation: ``Comercio'' Firms Win More on Diverse Tenders}
\textbf{Alternative test:} Use SII \textit{rubro} classification instead of bidding scope.\\
63\% of generalists are registered in \textit{Comercio al por mayor y menor} (wholesale/retail). \\[6pt]
\textbf{Outcome: Pr(Comercio firm wins the tender)}

\centering\small
\begin{tabular}{lcccc}
\toprule
 & \multicolumn{2}{c}{Pre-reform} & \multicolumn{2}{c}{Post-reform} \\
 & No seg & + Seg dum & No seg & + Seg dum \\
 & (1) & (2) & (3) & (4) \\
\midrule
N product codes & 0.0137$^{***}$ & 0.0047$^{***}$ & 0.0147$^{***}$ & 0.0064$^{***}$ \\
 & (0.0009) & (0.0007) & (0.0012) & (0.0010) \\[2pt]
N segments & 0.0372$^{***}$ & $-$0.1396$^{***}$ & 0.0321$^{***}$ & $-$0.1438$^{***}$ \\
 & (0.0025) & (0.0036) & (0.0030) & (0.0044) \\[2pt]
log(Est.\ cost) & $-$0.0376$^{***}$ & $-$0.0290$^{***}$ & $-$0.0408$^{***}$ & $-$0.0312$^{***}$ \\
 & (0.0022) & (0.0018) & (0.0021) & (0.0018) \\[2pt]
N bidders & 0.0249$^{***}$ & 0.0171$^{***}$ & 0.0210$^{***}$ & 0.0141$^{***}$ \\
 & (0.0005) & (0.0004) & (0.0005) & (0.0003) \\[2pt]
\midrule
Segment dummies & No & Yes & No & Yes \\
N & 726,056 & 726,056 & 219,511 & 219,511 \\
Within $R^2$ & 0.082 & 0.164 & 0.079 & 0.172 \\
\bottomrule
\multicolumn{5}{l}{\footnotesize Buyer + year-month FE. SEs clustered by buyer. $^{***}p<0.01$, $^{**}p<0.05$.} \\
\end{tabular}
\end{frame}

\begin{frame}{Comercio Wins: Adding Comercio Share Control}
\textbf{Decomposition:} Does the segment effect operate through selection (who bids) or competitive advantage (who wins conditional on bidding)?

\centering\small
\begin{tabular}{lcccc}
\toprule
 & \multicolumn{2}{c}{Pre-reform} & \multicolumn{2}{c}{Post-reform} \\
 & No seg & + Seg dum & No seg & + Seg dum \\
 & (1) & (2) & (3) & (4) \\
\midrule
N product codes & 0.0016$^{***}$ & 0.0010$^{**}$ & 0.0024$^{***}$ & 0.0013$^{**}$ \\
 & (0.0005) & (0.0005) & (0.0006) & (0.0006) \\[2pt]
N segments & 0.0085$^{***}$ & $-$0.0145$^{***}$ & 0.0059$^{***}$ & $-$0.0158$^{***}$ \\
 & (0.0012) & (0.0014) & (0.0016) & (0.0020) \\[2pt]
Share comercio bidders & 1.0088$^{***}$ & 0.9940$^{***}$ & 1.0005$^{***}$ & 0.9782$^{***}$ \\
 & (0.0012) & (0.0016) & (0.0021) & (0.0028) \\[2pt]
\midrule
Segment dummies & No & Yes & No & Yes \\
N & 726,056 & 726,056 & 219,511 & 219,511 \\
Within $R^2$ & 0.559 & 0.561 & 0.555 & 0.558 \\
\bottomrule
\multicolumn{5}{l}{\footnotesize Buyer + YM FE. Controls: log(cost), N bidders. SEs clustered by buyer.} \\
\end{tabular}

\bigskip
\small\textbf{Takeaway:} Comercio share of bidders is nearly a sufficient statistic ($R^2 \approx 0.56$). \\
Residual N\_segment2 effect is small: selection (who bids) dominates for comercio firms.
\end{frame}

\begin{frame}{Selection Channel: Share of Comercio Bidders}
\centering\small
\begin{tabular}{lcccc}
\toprule
 & \multicolumn{2}{c}{Pre-reform} & \multicolumn{2}{c}{Post-reform} \\
 & No seg & + Seg dum & No seg & + Seg dum \\
 & (1) & (2) & (3) & (4) \\
\midrule
N product codes & 0.0120$^{***}$ & 0.0037$^{***}$ & 0.0123$^{***}$ & 0.0052$^{***}$ \\
 & (0.0006) & (0.0004) & (0.0009) & (0.0007) \\[2pt]
N segments & 0.0284$^{***}$ & $-$0.1258$^{***}$ & 0.0262$^{***}$ & $-$0.1309$^{***}$ \\
 & (0.0017) & (0.0029) & (0.0022) & (0.0037) \\[2pt]
log(Est.\ cost) & $-$0.0290$^{***}$ & $-$0.0219$^{***}$ & $-$0.0376$^{***}$ & $-$0.0291$^{***}$ \\
 & (0.0019) & (0.0015) & (0.0016) & (0.0014) \\[2pt]
N bidders & 0.0211$^{***}$ & 0.0142$^{***}$ & 0.0185$^{***}$ & 0.0126$^{***}$ \\
 & (0.0004) & (0.0003) & (0.0004) & (0.0003) \\[2pt]
\midrule
Segment dummies & No & Yes & No & Yes \\
N & 726,056 & 726,056 & 219,511 & 219,511 \\
Within $R^2$ & 0.110 & 0.237 & 0.112 & 0.248 \\
\bottomrule
\multicolumn{5}{l}{\footnotesize Outcome: share of bidders classified as ``Comercio'' in SII.} \\
\multicolumn{5}{l}{\footnotesize Buyer + YM FE. SEs clustered by buyer. $^{***}p<0.01$.} \\
\end{tabular}

\bigskip
\small\textbf{Takeaway:} Without segment dummies, N\_segment2 increases the comercio share (+0.03$^{***}$). \\
With segment dummies, the sign flips ($-$0.13$^{***}$) --- the selection into comercio is driven by \textit{which} segments, not how many.
\end{frame}

\begin{frame}{Which Categories Favor Comercio Firms?}
\centering
\includegraphics[width=\textwidth]{output/product_mix/figures/segment_coef_comercio.pdf}

\smallskip
\small
Pharmaceuticals, medical supplies, and cleaning products show the largest comercio advantage. \\
Food \& beverages and construction materials show the smallest --- dominated by specialized producers.
\end{frame}

% ══════════════════════════════════════════════════════════
\section{Summary}
% ══════════════════════════════════════════════════════════

\begin{frame}{Putting It Together}
\textbf{Two dimensions of bundling have different effects:}

\bigskip
\begin{columns}[T]
  \column{0.48\textwidth}
  \textbf{Product variety} (N\_product8) \\[4pt]
  More distinct product codes overall
  \begin{itemize}
    \item Attracts large firms into the bidder pool
    \item Repels SME bidders
    \item Raises Pr(large firm wins)
    \item But: effect on \textit{winning} is entirely selection \\
          (disappears with SME share control)
  \end{itemize}

  \column{0.48\textwidth}
  \textbf{Category breadth} (N\_segment2) \\[4pt]
  More distinct 2-digit UNSPSC segments
  \begin{itemize}
    \item Shrinks total bidder pool
    \item But tilts composition toward local SMEs
    \item Raises Pr(local SME wins) --- even conditional on bidder composition
    \item Genuine competitive advantage for local generalists
  \end{itemize}
\end{columns}
\end{frame}

\begin{frame}{Implications for the Reform}
\begin{enumerate}
  \item \textbf{No fundamental tension} between bundling and SME preference:
  \begin{itemize}
    \item The apparent negative effect of product variety on SME outcomes operates through \textit{selection} (who bids), not competitive disadvantage
    \item Conditional on bidder composition, bundling does not hurt SMEs
  \end{itemize}

  \bigskip
  \item \textbf{Cross-category bundling complements local preference:}
  \begin{itemize}
    \item Tenders spanning multiple UNSPSC segments give local SMEs a genuine competitive edge
    \item The ``general store'' advantage: local generalists can source diverse products that distant specialists cannot
    \item This effect is stable pre and post reform, and robust to buyer FE, region FE, and segment controls
  \end{itemize}

  \bigskip
  \item \textbf{Policy implication:}
  \begin{itemize}
    \item The reform's induced bundling is not a threat to its SME goals
    \item If anything, cross-category consolidation amplifies the local SME advantage
  \end{itemize}
\end{enumerate}
\end{frame}

\begin{frame}{Contribution to the Literature}
\begin{enumerate}
  \item \textbf{Bundling as multi-dimensional treatment.}
        Most empirical work treats bundling as a scalar (bundle vs.\ unbundle, or contract size).
        Following Ridderstedt \& Nilsson (2023), we decompose into product variety and category breadth
        --- and show they have \textit{opposite} effects on SME outcomes.

  \bigskip
  \item \textbf{Selection vs.\ competitive advantage --- a clean decomposition.}
        The literature emphasizes entry deterrence (Estache \& Iimi 2011; Suzuki 2021)
        but rarely separates whether bundling hurts small firms through who \textit{bids}
        vs.\ who \textit{wins conditional on bidding}. We show the product variety channel
        is entirely selection; the category breadth channel is genuine competitive advantage.

  \bigskip
  \item \textbf{Bundling $\times$ preferential policies.}
        No prior work (to our knowledge) studies how procurement bundling interacts with
        SME/local preference policies. We show the interaction is \textit{complementary},
        not adversarial --- a policy-relevant result for the many countries combining
        simplified procurement with domestic preference.

  \bigskip
  \item \textbf{Firm-level mechanism: generalist advantage.}
        Bid-level within-tender FE regressions show the generalist penalty shrinks
        by $\sim$60\% post-reform and nearly vanishes on diverse tenders ---
        identifying the ``general store'' channel at the firm level.
\end{enumerate}
\end{frame}

% ══════════════════════════════════════════════════════════
\appendix
% ══════════════════════════════════════════════════════════


\begin{frame}[noframenumbering]{Robustness: Pr(Local SME Wins) --- Region FE}
\centering\small
\begin{tabular}{lcccc}
\toprule
 & Pre (Buyer) & Pre (Region) & Post (Buyer) & Post (Region) \\
 & (1) & (2) & (3) & (4) \\
\midrule
N product codes (8-digit) & 0.0039$^{***}$ & 0.0039$^{***}$ & -0.0006 & -0.0007 \\
 & (0.0007) & (0.0007) & (0.0008) & (0.0008) \\[2pt]
N segments (2-digit) & 0.0496$^{***}$ & 0.0507$^{***}$ & 0.0577$^{***}$ & 0.0583$^{***}$ \\
 & (0.0029) & (0.0030) & (0.0035) & (0.0037) \\[2pt]
Share SME bidders & 0.6594$^{***}$ & 0.6669$^{***}$ & 0.6317$^{***}$ & 0.6382$^{***}$ \\
 & (0.0148) & (0.0147) & (0.0173) & (0.0172) \\[2pt]
log(Estimated cost) & 0.0065$^{***}$ & 0.0058$^{***}$ & -0.0003 & -0.0006 \\
 & (0.0018) & (0.0020) & (0.0018) & (0.0022) \\[2pt]
N bidders & -0.0113$^{***}$ & -0.0115$^{***}$ & -0.0105$^{***}$ & -0.0107$^{***}$ \\
 & (0.0005) & (0.0005) & (0.0006) & (0.0007) \\[2pt]
\midrule
Fixed Effects & Buyer + YM & Region + YM & Buyer + YM & Region + YM \\
Segment dummies & Yes & Yes & Yes & Yes \\
SME share control & Yes & Yes & Yes & Yes \\
N & 545,518 & 545,518 & 155,433 & 155,433 \\
Within $R^2$ & 0.187 & 0.199 & 0.151 & 0.162 \\
\bottomrule
\multicolumn{5}{l}{\footnotesize SEs clustered by buyer. $^{***}p<0.01$, $^{**}p<0.05$, $^{*}p<0.10$.} \\
\end{tabular}
\end{frame}

\begin{frame}[noframenumbering]{Robustness: Share Local SME Bidders --- Region FE}
\centering\small
\begin{tabular}{lcccc}
\toprule
 & Pre (Buyer) & Pre (Region) & Post (Buyer) & Post (Region) \\
 & (1) & (2) & (3) & (4) \\
\midrule
N product codes (8-digit) & -0.0019$^{***}$ & -0.0020$^{***}$ & 0.0007 & 0.0006 \\
 & (0.0005) & (0.0005) & (0.0005) & (0.0006) \\[2pt]
N segments (2-digit) & 0.0398$^{***}$ & 0.0408$^{***}$ & 0.0343$^{***}$ & 0.0347$^{***}$ \\
 & (0.0018) & (0.0019) & (0.0022) & (0.0023) \\[2pt]
log(Estimated cost) & -0.0050$^{***}$ & -0.0044$^{***}$ & -0.0038$^{***}$ & -0.0030$^{*}$ \\
 & (0.0013) & (0.0015) & (0.0012) & (0.0015) \\[2pt]
N bidders & -0.0051$^{***}$ & -0.0051$^{***}$ & -0.0052$^{***}$ & -0.0053$^{***}$ \\
 & (0.0003) & (0.0003) & (0.0004) & (0.0004) \\[2pt]
\midrule
Fixed Effects & Buyer + YM & Region + YM & Buyer + YM & Region + YM \\
Segment dummies & Yes & Yes & Yes & Yes \\
N & 880,029 & 880,029 & 288,525 & 288,525 \\
Within $R^2$ & 0.110 & 0.119 & 0.109 & 0.118 \\
\bottomrule
\multicolumn{5}{l}{\footnotesize SEs clustered by buyer. $^{***}p<0.01$, $^{**}p<0.05$, $^{*}p<0.10$.} \\
\end{tabular}
\end{frame}

\begin{frame}[noframenumbering]{Robustness: N Total Bidders --- Region FE}
\centering\small
\begin{tabular}{lcccc}
\toprule
 & Pre (Buyer) & Pre (Region) & Post (Buyer) & Post (Region) \\
 & (1) & (2) & (3) & (4) \\
\midrule
N product codes (8-digit) & 0.0069 & 0.0163$^{***}$ & -0.0061 & 0.0002 \\
 & (0.0065) & (0.0063) & (0.0060) & (0.0061) \\[2pt]
N segments (2-digit) & -1.0823$^{***}$ & -1.1060$^{***}$ & -1.2045$^{***}$ & -1.2233$^{***}$ \\
 & (0.0236) & (0.0239) & (0.0300) & (0.0316) \\[2pt]
log(Estimated cost) & 0.3180$^{***}$ & 0.3039$^{***}$ & 0.2992$^{***}$ & 0.2637$^{***}$ \\
 & (0.0134) & (0.0149) & (0.0151) & (0.0174) \\[2pt]
\midrule
Fixed Effects & Buyer + YM & Region + YM & Buyer + YM & Region + YM \\
Segment dummies & Yes & Yes & Yes & Yes \\
N & 930,247 & 930,247 & 307,298 & 307,298 \\
Within $R^2$ & 0.120 & 0.124 & 0.111 & 0.111 \\
\bottomrule
\multicolumn{5}{l}{\footnotesize SEs clustered by buyer. $^{***}p<0.01$, $^{**}p<0.05$, $^{*}p<0.10$.} \\
\end{tabular}
\end{frame}

\end{document}
